% Session 20: ME Apps (last session of the MLCoEP series)
% Focus: agriculture, mechanical and manufacturing examples only. Generic demos are commented, not deleted
% (Optional: restore by uncommenting).
\section[ME Apps]{Session 20: ME Apps}
%%%%%%%%%%%%%%%%%%%%%%%%%%%%%%%%%%%%%%%%%%%%%%%%%%%%%%%%%%%%%%%%%%%%%%%%%%%%%%%%%%
\begin{frame}[fragile]\frametitle{}
\begin{center}
{\Large Applications to Mechanical Engineering}
\end{center}
\end{frame}

%%%%%%%%%%%%%%%%%%%%%%%%%%%%%%%%%%%%%%%%%%%%%%%%%%%%%%%%%%%
\begin{frame}[fragile]\frametitle{The Announcement}
\begin{center}
{\Large ``Revitalizing manufacturing through AI''}
\end{center}
\begin{itemize}
\item AI is already transforming the IT industry.
\item It is now time for an AI-powered society.
\item We will initially focus on the manufacturing industry.
\end{itemize}

\tiny{(Reference: ``Medium'' - Andrew Ng, Dec 14,2017)}
\end{frame}

%%%%%%%%%%%%%%%%%%%%%%%%%%%%%%%%%%%%%%%%%%%%%%%%%%%%%%%%%%%
\begin{frame}[fragile]\frametitle{Why Manufacturing?}
\begin{itemize}
\item Manufacturing touches every part of society.
\item Need to help not just bits and pixels but physical world.
\end{itemize}

\end{frame}

%%%%%%%%%%%%%%%%%%%%%%%%%%%%%%%%%%%%%%%%%%%%%%%%%%%%%%%%%%%
\begin{frame}[fragile]\frametitle{Challenges Facing Manufacturing}
\begin{itemize}
\item Variable quality and yield,
\item Inflexible production line design
\item Inability to manage capacity
\item and rising production costs.
\end{itemize}
ML to help address them.
\end{frame}

%%%%%%%%%%%%%%%%%%%%%%%%%%%%%%%%%%%%%%%%%%%%%%%%%%%%%%%%%%%
\begin{frame}[fragile]\frametitle{Challenges Facing Manufacturing: More Ways ML Helps}
\begin{itemize}
\item Shorten design cycles,
\item Remove supply-chain bottlenecks,
\item Reduce materials and energy waste,
\item and improve production yields.
\end{itemize}

\end{frame}


%%%%%%%%%%%%%%%%%%%%%%%%%%%%%%%%%%%%%%%%%%%%%%%%%%%%%%%%%%%
\begin{frame}[fragile]\frametitle{Why now?}
\begin{center}
{\Large ``Industry today is experiencing a never seen increase in available data'' }
\end{center}
Sources:
\begin{itemize}
\item Environmental data
\item Sensor data from production line
\item Machine tool parameters
\item Logistics, transaction data
\item Customer demand data.
\end{itemize}
\tiny{(Reference: What is smart manufacturing? Time Magazine - Chand \& Davis, 2010)}
\end{frame}

%%%%%%%%%%%%%%%%%%%%%%%%%%%%%%%%%%%%%%%%%%%%%%%%%%%%%%%%%%%
\begin{frame}[fragile]\frametitle{How ML can help?: High-Dimensional Data}

\begin{itemize}
\item Ability to handle high-dimensional problems
\item SVM handles high dimensionality ($>1000$ features) very well.
\item Concerns: Over-fitting
\end{itemize}
{\tiny (Ref: ``Machine learning in manufacturing: advantages, challenges, and applications''- Thorsten Wuest, et al)}
\end{frame}


%%%%%%%%%%%%%%%%%%%%%%%%%%%%%%%%%%%%%%%%%%%%%%%%%%%%%%%%%%%
\begin{frame}[fragile]\frametitle{How ML can help?: Interpretability}

\begin{itemize}
\item Ability to reduce complexity
\item Present transparent and concrete advice for practitioners
\item E.g. monitor XX and parameter YY at checkpoint ZZ
\item ML finds important features, thresholds
\end{itemize}
{\tiny (Ref: ``Machine learning in manufacturing: advantages, challenges, and applications''- Thorsten Wuest, et al)}
\end{frame}

%%%%%%%%%%%%%%%%%%%%%%%%%%%%%%%%%%%%%%%%%%%%%%%%%%%%%%%%%%%
\begin{frame}[fragile]\frametitle{How ML can help?: Adaptability}

\begin{itemize}
\item Ability to adapt to changing environment
\item New data, new training, new model.
\end{itemize}
{\tiny (Ref: ``Machine learning in manufacturing: advantages, challenges, and applications''- Thorsten Wuest, et al)}
\end{frame}

%%%%%%%%%%%%%%%%%%%%%%%%%%%%%%%%%%%%%%%%%%%%%%%%%%%%%%%%%%%
\begin{frame}[fragile]\frametitle{New Wave: What is IoT?}
``Internet of Things''
\begin{itemize}
\item IoT is connecting things to the Internet
\item Not the same as connecting things to the cellular network!
\item IoT personal (Health apps)
\item IoT Machine (Machine 2 Machine M2M)
\item IoT Everywhere (5G, wifi)
\end{itemize}
\end{frame}

%%%%%%%%%%%%%%%%%%%%%%%%%%%%%%%%%%%%%%%%%%%%%%%%%%%%%%%%%%%
\begin{frame}[fragile]\frametitle{New Wave: IoT Data}
\begin{itemize}
\item Sensors everywhere.
\item Tens of billions of connected devices worldwide, and still growing.
\item Wearables, cars, machines, cities.
\end{itemize}
\end{frame}

%%%%%%%%%%%%%%%%%%%%%%%%%%%%%%%%%%%%%%%%%%%%%%%%%%%%%%%%%%%
\begin{frame}[fragile]\frametitle{New Wave: IoT and Machine Learning}
\begin{itemize}
\item Build a mathematical model from training data (learning stage).
\item Predict results for new data (prediction stage).
\item Tweak the model based on new conditions.
\item IoT creates a lot of contextual data
\item Use the scale of IoT data to gain new insights
\end{itemize}
\end{frame}

%%%%%%%%%%%%%%%%%%%%%%%%%%%%%%%%%%%%%%%%%%%%%%%%%%%%%%%%%%%
\begin{frame}[fragile]\frametitle{New Wave: IoT ML Applications}
\begin{itemize}
\item Consumer: bio sensors (real time tracking)
\item Enterprise: complex machinery (preventative maintenance), asset efficiency.
\item Public infrastructure: Smart Cities, dynamically adjust traffic lights
\end{itemize}
\end{frame}
%%%%%%%%%%%%%%%%%%%%%%%%%%%%%%%%%%%%%%%%%%%%%%%%%%%%%%%%%%%
\begin{frame}[fragile]\frametitle{Sample Applications: Tool Wear Estimation}
 Hidden Markov Model
\begin{itemize}
\item Hidden states: the wear stage of the tool (for example sharp, worn, failed)
\item Observations: sensor signals (for example cutting force or vibration)
\item The model infers the most likely wear stage from the signals
\end{itemize}
\begin{center}
\includegraphics[width=0.55\linewidth,keepaspectratio]{mechapp2}
\end{center}
\end{frame}



% %%%%%%%%%%%%%%%%%%%%%%%%%%%%%%%%%%%%%%%%%%%%%%%%%%%%%%%%%%%
% \begin{frame}[fragile]\frametitle{Other Applications: Robot  Playing Piano}

% \begin{center}
% \includegraphics[width=0.8\linewidth,keepaspectratio]{mechapp1}
% \end{center}
% {\tiny (Ref: Machine Learning and Deep Learning in Manufacturing - Prof. Seungchul Lee)}
% \end{frame}

%%%%%%%%%%%%%%%%%%%%%%%%%%%%%%%%%%%%%%%%%%%%%%%%%%%%%%%%%%%
\begin{frame}[fragile]\frametitle{Ideas for new applications}
Project ideas?
\begin{itemize}
\item Kaggle: \url{https://www.kaggle.com/datasets}
\item UCI: \url{https://archive.ics.uci.edu}
\item Hugging Face: \url{https://huggingface.co/datasets}
\item NASA Prognostics Data Repository: turbofan engine and bearing run-to-failure data (search for the name)
\end{itemize}
\end{frame}

%%%%%%%%%%%%%%%%%%%%%%%%%%%%%%%%%%%%%%%%%%%%%%%%%%%%%%%%%%%%
\begin{frame}[fragile]\frametitle{Quick Check: ML in Manufacturing}
\begin{block}{Think About It}
A manufacturing line records hundreds of sensor readings per part (temperature, vibration,
tool position, timing, etc.) and needs a model to predict good vs. defective output. Why
was SVM specifically called out as a good algorithmic fit for this kind of problem?
\end{block}
\end{frame}

%%%%%%%%%%%%%%%%%%%%%%%%%%%%%%%%%%%%%%%%%%%%%%%%%%%%%%%%%%%%
\begin{frame}[fragile]\frametitle{Quick Check: ML in Manufacturing}
\begin{block}{Answer}
SVM handles high-dimensional data (hundreds or thousands of features) well, and sensor-rich
manufacturing lines are exactly this kind of high-dimensional problem: each sensor reading
is a feature, so a part can easily have hundreds of them. The main concern flagged along
with it is over-fitting, since more features relative to the number of training examples
makes it easier for the model to memorize noise instead of learning the real pattern.
\end{block}
\end{frame}






%%%%%%%%%%%%%%%%%%%%%%%%%%%%%%%%%%%%%%%%%%%%%%%%%%%%%%%%%%%%
%\begin{frame}[fragile]\frametitle{Some types of algorithms}
%\begin{itemize}
%\item Prediction: predicting a variable from data
%\item Classification: assigning records to predefined groups
%\item Clustering: splitting records into groups based on similarity
%\item Association learning: seeing what often appears together
%\end{itemize}
%
%\end{frame}
%
%
%
%%%%%%%%%%%%%%%%%%%%%%%%%%%%%%%%%%%%%%%%%%%%%%%%%%%%%%%%%%%%
%\begin{frame}[fragile]\frametitle{How to figure out? A dumb way}
%Step 1: Start with each weight set to 1.0:
%\begin{lstlisting}
%def estimate_house_sales_price(num_of_bedrooms, sqft, neighborhood):
%	price = 0
%	price += num_of_bedrooms * 1.0
%	price += sqft * 1.0
%	price += neighborhood * 1.0
%	price += 1.0
%
% return price
%\end{lstlisting}
%\end{frame}
%




% Optional (off-theme: not agriculture, mechanical or manufacturing). Uncomment to restore.
% %%%%%%%%%%%%%%%%%%%%%%%%%%%%%%%%%%%%%%%%%%%%%%%%%%%%%%%%%%%%%%%%%%%%%%%%%%%%%%%%%%
\begin{frame}[fragile]\frametitle{}
\begin{center}
{\Large Linear Regression: House Price Prediction}

{\tiny (Ref: scikit-learn User Guide: Generalized Linear Models)}
\end{center}
\end{frame}

%%%%%%%%%%%%%%%%%%%%%%%%%%%%%%%%%%%%%%%%%%%%%%%%%%%%%%%%%%
\begin{frame}[fragile]\frametitle{Problem Info}
\begin{itemize}
\item Dataset: California Housing (\texttt{sklearn.\allowbreak datasets.\allowbreak fetch\_california\_housing}, downloaded on first use)
\item Features: median income, house age, avg rooms, avg bedrooms, population, avg occupancy, latitude, longitude
\item Target: Median house value (in 100,000s USD)
\item Goal: Predict house prices using Linear Regression; compare Ridge and Lasso regularization
\end{itemize}
\end{frame}

%%%%%%%%%%%%%%%%%%%%%%%%%%%%%%%%%%%%%%%%%%%%%%%%%%%%%%%%%%
\begin{frame}[fragile]\frametitle{Import Libraries}
\begin{lstlisting}
import numpy as np
import pandas as pd
import matplotlib.pyplot as plt
from sklearn.datasets import fetch_california_housing
from sklearn.model_selection import train_test_split
from sklearn.linear_model import LinearRegression, Ridge, Lasso
from sklearn.metrics import mean_squared_error, r2_score
from sklearn.preprocessing import StandardScaler
\end{lstlisting}
\end{frame}

%%%%%%%%%%%%%%%%%%%%%%%%%%%%%%%%%%%%%%%%%%%%%%%%%%%%%%%%%%
\begin{frame}[fragile]\frametitle{Load and Explore Data}
\begin{lstlisting}
data = fetch_california_housing()
df = pd.DataFrame(data.data, columns=data.feature_names)
df['MedHouseVal'] = data.target

print(df.shape)       # (20640, 9)
print(df.describe())
print(df.head())
\end{lstlisting}
\end{frame}

%%%%%%%%%%%%%%%%%%%%%%%%%%%%%%%%%%%%%%%%%%%%%%%%%%%%%%%%%%
\begin{frame}[fragile]\frametitle{Train Linear Regression}
\begin{itemize}
\item RMSE: the typical size of the prediction error, in the units of the target
\item R2: the share of the variation in prices that the model explains (1 is perfect)
\end{itemize}
\begin{lstlisting}
X_train, X_test, y_train, y_test = train_test_split(
    data.data, data.target, test_size=0.2, random_state=42)

scaler = StandardScaler()
X_train = scaler.fit_transform(X_train)
X_test  = scaler.transform(X_test)

lr = LinearRegression()
lr.fit(X_train, y_train)
y_pred = lr.predict(X_test)

print("RMSE:", np.sqrt(mean_squared_error(y_test, y_pred)))
print("R2  :", r2_score(y_test, y_pred))
\end{lstlisting}
\end{frame}

%%%%%%%%%%%%%%%%%%%%%%%%%%%%%%%%%%%%%%%%%%%%%%%%%%%%%%%%%%
\begin{frame}[fragile]\frametitle{Compare Regularization}
\begin{itemize}
\item Regularization adds a penalty on large coefficients, to keep the model simple
\item Ridge shrinks all coefficients; Lasso can shrink some to exactly 0; \texttt{alpha} sets the strength of the penalty
\end{itemize}
\begin{lstlisting}
for name, model in [("Linear", LinearRegression()),
                    ("Ridge",  Ridge(alpha=1.0)),
                    ("Lasso",  Lasso(alpha=0.1))]:
    model.fit(X_train, y_train)
    pred = model.predict(X_test)
    rmse = np.sqrt(mean_squared_error(y_test, pred))
    r2   = r2_score(y_test, pred)
    print(f"{name:8s}  RMSE={rmse:.4f}  R2={r2:.4f}")
\end{lstlisting}
\end{frame}

%%%%%%%%%%%%%%%%%%%%%%%%%%%%%%%%%%%%%%%%%%%%%%%%%%%%%%%%%%
\begin{frame}[fragile]\frametitle{Predicted vs Actual}
Points close to the diagonal indicate accurate predictions.
\begin{lstlisting}
plt.figure(figsize=(6, 5))
plt.scatter(y_test, y_pred, alpha=0.3, s=10)
plt.plot([y_test.min(), y_test.max()],
         [y_test.min(), y_test.max()], 'r--', lw=2)
plt.xlabel("Actual Price")
plt.ylabel("Predicted Price")
plt.title("Linear Regression: Predicted vs Actual")
plt.tight_layout()
plt.show()
\end{lstlisting}
\end{frame}

%%%%%%%%%%%%%%%%%%%%%%%%%%%%%%%%%%%%%%%%%%%%%%%%%%%
\begin{frame}[fragile]\frametitle{Results: Regularization Comparison}
\begin{itemize}
\item Linear and Ridge score the same: with 8 standardized features and 16,512 training rows, regularization adds little
\item Lasso with $\alpha = 0.1$ penalizes too strongly here, and its scores are lower
\item RMSE 0.75 means a typical error of about \$75,000, since the target is in units of \$100,000
\end{itemize}
\begin{lstlisting}
Linear    RMSE=0.7456  R2=0.5758
Ridge     RMSE=0.7456  R2=0.5758
Lasso     RMSE=0.8244  R2=0.4814
\end{lstlisting}
\end{frame}

%%%%%%%%%%%%%%%%%%%%%%%%%%%%%%%%%%%%%%%%%%%%%%%%%%%
\begin{frame}[fragile]\frametitle{Quick Check: Ridge and Lasso}
\begin{block}{Think About It}
Linear regression and Ridge scored the same on the housing data, and Lasso with $\alpha = 0.1$ scored lower. What is the best explanation?
\begin{enumerate}
  \item[A.] With 8 useful features and thousands of rows the plain model barely overfits, so the penalty only removes useful signal
  \item[B.] Lasso cannot be used on standardized features
  \item[C.] Ridge always beats Lasso
  \item[D.] Regularization always improves a model, so the result must be a bug
\end{enumerate}
\end{block}
\end{frame}

%%%%%%%%%%%%%%%%%%%%%%%%%%%%%%%%%%%%%%%%%%%%%%%%%%%
\begin{frame}[fragile]\frametitle{Quick Check: Ridge and Lasso}
\begin{block}{Answer}
\textbf{Correct: A.} Regularization helps when a model overfits: many features, few rows, or noisy data. Here there is little overfitting to remove, so the penalty only shrinks coefficients that carry real signal, and a strong penalty ($\alpha = 0.1$ for Lasso) costs accuracy. Lasso works on standardized features (B), neither method always wins (C), and regularization is a tool to tune, not a guaranteed gain (D).
\end{block}
\end{frame}
% %%%%%%%%%%%%%%%%%%%%%%%%%%%%%%%%%%%%%%%%%%%%%%%%%%%%%%%%%%%%%%%%%%%%%%%%%%%%%%%%%%
\begin{frame}[fragile]\frametitle{}
\begin{center}
{\Large SVM: Handwritten Digit Recognition}

{\tiny (Ref: scikit-learn SVM example: Digit Classification)}
\end{center}
\end{frame}

%%%%%%%%%%%%%%%%%%%%%%%%%%%%%%%%%%%%%%%%%%%%%%%%%%%%%%%%%%
\begin{frame}[fragile]\frametitle{Problem Info}
\begin{itemize}
\item Dataset: Digits (\texttt{sklearn.datasets.load\_digits}, built-in)
\item 1,797 images of handwritten digits (0--9), each 8$\times$8 pixels = 64 features
\item Goal: Classify each image into one of 10 digit classes using SVM
\item Demonstrates the power of SVM's RBF kernel for non-linear image classification
\end{itemize}
\end{frame}

%%%%%%%%%%%%%%%%%%%%%%%%%%%%%%%%%%%%%%%%%%%%%%%%%%%%%%%%%%
\begin{frame}[fragile]\frametitle{Import Libraries}
\begin{lstlisting}
import numpy as np
import matplotlib.pyplot as plt
from sklearn.datasets import load_digits
from sklearn.model_selection import train_test_split
from sklearn.svm import SVC
from sklearn.preprocessing import StandardScaler
from sklearn.metrics import classification_report, confusion_matrix
import seaborn as sns
\end{lstlisting}
\end{frame}

%%%%%%%%%%%%%%%%%%%%%%%%%%%%%%%%%%%%%%%%%%%%%%%%%%%%%%%%%%
\begin{frame}[fragile]\frametitle{Load and Visualize Data}
\begin{lstlisting}
digits = load_digits()
print("Data shape:", digits.data.shape)   # (1797, 64)
print("Classes   :", digits.target_names) # [0 1 2 ... 9]

fig, axes = plt.subplots(2, 5, figsize=(10, 4))
for ax, img, label in zip(axes.ravel(),
                           digits.images, digits.target):
    ax.imshow(img, cmap='gray_r')
    ax.set_title(str(label))
    ax.axis('off')
plt.suptitle("Sample Digits")
plt.show()
\end{lstlisting}
\end{frame}

%%%%%%%%%%%%%%%%%%%%%%%%%%%%%%%%%%%%%%%%%%%%%%%%%%%%%%%%%%
\begin{frame}[fragile]\frametitle{Train SVM Classifier}
\begin{itemize}
\item \texttt{C}: how strictly the model fits the training points (large C: fewer training errors, higher risk of overfitting)
\item \texttt{gamma}: how far the influence of a single training point reaches (small gamma: far)
\end{itemize}
\begin{lstlisting}
X_train, X_test, y_train, y_test = train_test_split(
    digits.data, digits.target,
    test_size=0.2, random_state=42)

scaler = StandardScaler()
X_train_scaled = scaler.fit_transform(X_train)
X_test_scaled  = scaler.transform(X_test)

svm = SVC(kernel='rbf', C=10, gamma=0.001)
svm.fit(X_train_scaled, y_train)
y_pred = svm.predict(X_test_scaled)

print(classification_report(y_test, y_pred))
\end{lstlisting}
\end{frame}

%%%%%%%%%%%%%%%%%%%%%%%%%%%%%%%%%%%%%%%%%%%%%%%%%%%%%%%%%%
\begin{frame}[fragile]\frametitle{Confusion Matrix}
\begin{lstlisting}
cm = confusion_matrix(y_test, y_pred)
plt.figure(figsize=(8, 6))
sns.heatmap(cm, annot=True, fmt='d', cmap='Blues',
            xticklabels=digits.target_names,
            yticklabels=digits.target_names)
plt.xlabel("Predicted")
plt.ylabel("True")
plt.title("SVM Digit Recognition: Confusion Matrix")
plt.tight_layout()
plt.show()
\end{lstlisting}
\end{frame}

%%%%%%%%%%%%%%%%%%%%%%%%%%%%%%%%%%%%%%%%%%%%%%%%%%%%%%%%%%
\begin{frame}[fragile]\frametitle{Visualize Misclassifications}
\begin{lstlisting}
wrong = np.where(y_pred != y_test)[0]
fig, axes = plt.subplots(2, 5, figsize=(10, 4))
for ax, idx in zip(axes.ravel(), wrong[:10]):
    ax.imshow(X_test[idx].reshape(8, 8), cmap='gray_r')
    ax.set_title(f"True:{y_test[idx]} Pred:{y_pred[idx]}")
    ax.axis('off')
plt.suptitle("Misclassified Digits")
plt.tight_layout()
plt.show()
\end{lstlisting}
\end{frame}

%%%%%%%%%%%%%%%%%%%%%%%%%%%%%%%%%%%%%%%%%%%%%%%%%%%
\begin{frame}[fragile]\frametitle{Results: Test Accuracy}
\begin{itemize}
\item 98.1\% of the 360 test images are classified correctly (7 errors)
\item The RBF kernel on standardized pixels separates the digit classes well
\item The misclassified digits are often ones a person might also find ambiguous
\end{itemize}
\begin{lstlisting}
# accuracy on the 360 test images
0.9806
\end{lstlisting}
\end{frame}

%%%%%%%%%%%%%%%%%%%%%%%%%%%%%%%%%%%%%%%%%%%%%%%%%%%
\begin{frame}[fragile]\frametitle{Quick Check: The gamma Parameter}
\begin{block}{Think About It}
The SVM used \texttt{gamma=0.001} and scored 98\% on the test images. If you raise \texttt{gamma} to 10 and refit, what most likely happens?
\begin{enumerate}
  \item[A.] Test accuracy rises, because each training point now has more influence
  \item[B.] Training accuracy stays near 100\% but test accuracy collapses: the model memorizes the training images
  \item[C.] Nothing changes, because the data is already standardized
  \item[D.] The model underfits, because a large gamma makes the boundary too smooth
\end{enumerate}
\end{block}
\end{frame}

%%%%%%%%%%%%%%%%%%%%%%%%%%%%%%%%%%%%%%%%%%%%%%%%%%%
\begin{frame}[fragile]\frametitle{Quick Check: The gamma Parameter}
\begin{block}{Answer}
\textbf{Correct: B.} A large \texttt{gamma} makes the influence of each training point very narrow, so the boundary wraps around individual points. In a run with \texttt{gamma=10}, the training accuracy was 1.00 while the test accuracy fell to about 0.08. A small gamma gives a smooth boundary (D describes the opposite effect), and gamma matters even for standardized data (C).
\end{block}
\end{frame}
% %%%%%%%%%%%%%%%%%%%%%%%%%%%%%%%%%%%%%%%%%%%%%%%%%%%%%%%%%%%%%%%%%%%%%%%%%%%%%%%%%%
\begin{frame}[fragile]\frametitle{}
\begin{center}
{\Large K-Means Clustering: Customer Segmentation}

{\tiny (Ref: scikit-learn Clustering User Guide)}
\end{center}
\end{frame}

%%%%%%%%%%%%%%%%%%%%%%%%%%%%%%%%%%%%%%%%%%%%%%%%%%%%%%%%%%
\begin{frame}[fragile]\frametitle{Problem Info}
\begin{itemize}
\item Scenario: A retail company wants to segment its customers for targeted marketing
\item Features: Annual Income (k USD) and Spending Score (1--100)
\item Data: Synthetic dataset mimicking a retail loyalty program
\item Goal: Discover natural customer groups using K-Means; choose optimal K with the Elbow method
\end{itemize}
\end{frame}

%%%%%%%%%%%%%%%%%%%%%%%%%%%%%%%%%%%%%%%%%%%%%%%%%%%%%%%%%%
\begin{frame}[fragile]\frametitle{Import Libraries and Generate Data}
\begin{lstlisting}
import numpy as np
import pandas as pd
import matplotlib.pyplot as plt
from sklearn.cluster import KMeans
from sklearn.preprocessing import StandardScaler

np.random.seed(42)
# Simulate 3 customer segments: budget, standard, premium
centers = [(30, 20), (55, 55), (80, 85)]
X = np.vstack([
    np.random.randn(120, 2) * [8, 10] + c
    for c in centers])
df = pd.DataFrame(X, columns=['AnnualIncome', 'SpendingScore'])
print(df.describe())
\end{lstlisting}
\end{frame}

%%%%%%%%%%%%%%%%%%%%%%%%%%%%%%%%%%%%%%%%%%%%%%%%%%%%%%%%%%
\begin{frame}[fragile]\frametitle{Elbow Method: Choosing K}
\begin{itemize}
\item Inertia (WCSS, within-cluster sum of squares): the total squared distance of the points to their cluster centers
\item Look for the elbow: the K after which extra clusters help little
\end{itemize}
\begin{lstlisting}
scaler = StandardScaler()
X_scaled = scaler.fit_transform(df)

inertia = []
K_range = range(1, 10)
for k in K_range:
    km = KMeans(n_clusters=k, random_state=42, n_init=10)
    km.fit(X_scaled)
    inertia.append(km.inertia_)

plt.plot(K_range, inertia, marker='o')
plt.xlabel("Number of Clusters K")
plt.ylabel("Inertia (WCSS)")
plt.title("Elbow Method")
plt.grid(True); plt.show()
\end{lstlisting}
\end{frame}

%%%%%%%%%%%%%%%%%%%%%%%%%%%%%%%%%%%%%%%%%%%%%%%%%%%%%%%%%%
\begin{frame}[fragile]\frametitle{Fit K-Means with Optimal K}
\begin{lstlisting}
km = KMeans(n_clusters=3, random_state=42, n_init=10)
df['Segment'] = km.fit_predict(X_scaled)

# KMeans cluster numbers are arbitrary, so map them to
# Budget/Standard/Premium by sorting on mean income
order = df.groupby('Segment')['AnnualIncome'].mean().sort_values().index
labels = dict(zip(order, ['Budget', 'Standard', 'Premium']))
colors = {'Budget': '#e07b54', 'Standard': '#5b9bd5', 'Premium': '#70ad47'}

plt.figure(figsize=(7, 5))
for seg, name in labels.items():
    mask = df['Segment'] == seg
    plt.scatter(df.loc[mask, 'AnnualIncome'],
                df.loc[mask, 'SpendingScore'],
                label=name, color=colors[name], alpha=0.7)
plt.xlabel("Annual Income (k USD)")
plt.ylabel("Spending Score")
plt.title("Customer Segments (K-Means, K=3)")
plt.legend(); plt.tight_layout(); plt.show()
\end{lstlisting}
\end{frame}

%%%%%%%%%%%%%%%%%%%%%%%%%%%%%%%%%%%%%%%%%%%%%%%%%%%%%%%%%%
\begin{frame}[fragile]\frametitle{Fit K-Means with Optimal K: Intuition}
\begin{block}{Intuition}
KMeans cluster numbers (0, 1, 2, \ldots) are arbitrary and can come out in any order;
they carry no meaning by themselves. To turn them into business labels like ``Budget'' or
``Premium'', sort the clusters by a real statistic (here, mean income) after fitting,
rather than assuming cluster 0 is always the first segment.
\end{block}
\end{frame}

%%%%%%%%%%%%%%%%%%%%%%%%%%%%%%%%%%%%%%%%%%%%%%%%%%%%%%%%%%
\begin{frame}[fragile]\frametitle{Segment Profiles}
\begin{itemize}
\item \textbf{Budget}: low income, low spending; price-sensitive customers
\item \textbf{Standard}: mid income, mid spending; core customer base
\item \textbf{Premium}: high income, high spending; loyalty program targets
\end{itemize}
\begin{lstlisting}
profile = df.groupby('Segment')[
    ['AnnualIncome','SpendingScore']].mean()
profile.index = [labels[i] for i in profile.index]
print(profile.sort_values('AnnualIncome').round(1))
\end{lstlisting}
\end{frame}

%%%%%%%%%%%%%%%%%%%%%%%%%%%%%%%%%%%%%%%%%%%%%%%%%%%
\begin{frame}[fragile]\frametitle{Results: Elbow and Segments}
\begin{itemize}
\item Inertia drops sharply up to $K = 3$ (720, 217, 88) and then only slowly: the elbow is at $K = 3$
\item The three segments match the three groups that generated the data
\end{itemize}
\begin{lstlisting}
# inertia for K = 1, 2, 3, ...
720.0  216.9  87.7  76.5  65.2  54.8  48.2  41.4  37.3

# segment profiles: mean income, mean spending score
          AnnualIncome  SpendingScore
Budget            29.3           20.7
Standard          55.4           54.7
Premium           79.0           84.6
\end{lstlisting}
\end{frame}

% %%%%%%%%%%%%%%%%%%%%%%%%%%%%%%%%%%%%%%%%%%%%%%%%%%%%%%%%%%%%%%%%%%%%%%%%%%%%%%%%%%
\begin{frame}[fragile]\frametitle{}
\begin{center}
{\Large PCA + K-Means: Handwritten Digit Clustering}

{\tiny (Ref: scikit-learn Decomposition and Clustering User Guide)}
\end{center}
\end{frame}

%%%%%%%%%%%%%%%%%%%%%%%%%%%%%%%%%%%%%%%%%%%%%%%%%%%%%%%%%%
\begin{frame}[fragile]\frametitle{Problem Info}
\begin{itemize}
\item Dataset: Digits (\texttt{sklearn.datasets.load\_digits}, built-in); 1,797 images, 64 features each
\item Part A: Use PCA to reduce 64 dimensions to 2; visualize digit clusters in 2D
\item Part B: Apply K-Means (K=10) on PCA-reduced data; measure cluster quality
\item Part C: Vary the number of PCA components (2, 10, 30) and observe effect on clustering
\end{itemize}
\end{frame}

%%%%%%%%%%%%%%%%%%%%%%%%%%%%%%%%%%%%%%%%%%%%%%%%%%%%%%%%%%
\begin{frame}[fragile]\frametitle{Import Libraries}
\begin{lstlisting}
import numpy as np
import matplotlib.pyplot as plt
from sklearn.datasets import load_digits
from sklearn.preprocessing import StandardScaler
from sklearn.decomposition import PCA
from sklearn.cluster import KMeans
from sklearn.metrics import silhouette_score, adjusted_rand_score
\end{lstlisting}
\end{frame}

%%%%%%%%%%%%%%%%%%%%%%%%%%%%%%%%%%%%%%%%%%%%%%%%%%%%%%%%%%
\begin{frame}[fragile]\frametitle{Part A: PCA Visualization}
\begin{lstlisting}
digits = load_digits()
X = StandardScaler().fit_transform(digits.data)

pca2 = PCA(n_components=2)
X2   = pca2.fit_transform(X)

plt.figure(figsize=(8, 6))
scatter = plt.scatter(X2[:, 0], X2[:, 1],
                      c=digits.target, cmap='tab10',
                      s=10, alpha=0.7)
plt.colorbar(scatter, label='Digit')
plt.title("Digits projected to 2 PCA components")
plt.xlabel("PC 1"); plt.ylabel("PC 2")
plt.tight_layout(); plt.show()

print("Variance explained:", pca2.explained_variance_ratio_.sum())
\end{lstlisting}
\end{frame}

%%%%%%%%%%%%%%%%%%%%%%%%%%%%%%%%%%%%%%%%%%%%%%%%%%%%%%%%%%
\begin{frame}[fragile]\frametitle{Parts B \& C: K-Means with Varying PCA Components}
\begin{itemize}
\item \textbf{Silhouette score}: measures cluster compactness (higher is better, max 1.0)
\item \textbf{Adjusted Rand Index}: compares clusters to true digit labels (1.0 = perfect)
\end{itemize}
\begin{lstlisting}
for n_comp in [2, 10, 30]:
    Xr = PCA(n_components=n_comp).fit_transform(X)
    km = KMeans(n_clusters=10, random_state=42, n_init=10)
    labels = km.fit_predict(Xr)
    sil = silhouette_score(Xr, labels)
    ari = adjusted_rand_score(digits.target, labels)
    print(f"PCA={n_comp:2d}  Silhouette={sil:.3f}  "
          f"ARI={ari:.3f}")
\end{lstlisting}
\end{frame}

%%%%%%%%%%%%%%%%%%%%%%%%%%%%%%%%%%%%%%%%%%%%%%%%%%%%%%%%%%
\begin{frame}[fragile]\frametitle{Assignment Tasks}
\begin{itemize}
\item Report the variance explained by 2, 10, and 30 PCA components
\item Plot the explained variance ratio (scree plot) and identify the ``elbow''
\item At which number of PCA components does ARI peak? Explain why
\item Visualize cluster centers for the best configuration: keep the fitted PCA object, and undo both the PCA and the scaling to return to pixel space
\item What do the cluster centers look like? Do they resemble recognisable digits?
\end{itemize}
\end{frame}


%%%%%%%%%%%%%%%%%%%%%%%%%%%%%%%%%%%%%%%%%%%%%%%%%%%%%%%%%%%
\begin{frame}[fragile]\frametitle{}
\begin{center}
{\Large Assignments for Mechanical Engineering}
\end{center}
\end{frame}

%%%%%%%%%%%%%%%%%%%%%%%%%%%%%%%%%%%%%%%%%%%%%%%%%%%%%%%%%%%
% Slide 1: Predictive Maintenance using Regression Analysis
%%%%%%%%%%%%%%%%%%%%%%%%%%%%%%%%%%%%%%%%%%%%%%%%%%%%%%%%%%%
\begin{frame}[fragile]\frametitle{Predictive Maintenance using Regression Analysis}
    \begin{itemize}
        \item \textbf{Dataset}: NASA Turbofan Engine Degradation Simulation Data Set
        \item \textbf{Source}: \href{https://www.nasa.gov/content/prognostics-center-of-excellence-data-set-repository}{NASA Prognostics Data Repository}
        \item \textbf{Description}: Predict the remaining useful life of turbofan engines using regression analysis based on sensor data.
    \end{itemize}
    \begin{lstlisting}[language=Python]
import pandas as pd
from sklearn.model_selection import train_test_split
from sklearn.linear_model import LinearRegression

data = pd.read_csv('turbofan_data.csv')
X = data[['feature1', 'feature2', 'feature3']]
y = data['remaining_useful_life']

X_train, X_test, y_train, y_test = train_test_split(X, y, test_size=0.2)
model = LinearRegression()
model.fit(X_train, y_train)
predictions = model.predict(X_test)
    \end{lstlisting}
\end{frame}

%%%%%%%%%%%%%%%%%%%%%%%%%%%%%%%%%%%%%%%%%%%%%%%%%%%%%%%%%%%
% Slide 2: Structural Health Monitoring using Anomaly Detection
%%%%%%%%%%%%%%%%%%%%%%%%%%%%%%%%%%%%%%%%%%%%%%%%%%%%%%%%%%%
\begin{frame}[fragile]\frametitle{Structural Health Monitoring using Anomaly Detection}
    \begin{itemize}
        \item \textbf{Dataset}: A public structural health monitoring dataset (vibration or strain sensors)
        \item \textbf{Source}: \href{https://www.kaggle.com/datasets?search=structural+health+monitoring}{Kaggle search}
        \item \textbf{Description}: Detect anomalies indicating potential structural failures in bridges based on sensor readings.
    \end{itemize}
    \begin{lstlisting}[language=Python]
import pandas as pd
from sklearn.ensemble import IsolationForest

data = pd.read_csv('bridge_health_data.csv')
model = IsolationForest(contamination=0.1)   # expected share of anomalies
model.fit(data)
anomalies = model.predict(data)     # -1 = anomaly, 1 = normal
    \end{lstlisting}
\end{frame}

%%%%%%%%%%%%%%%%%%%%%%%%%%%%%%%%%%%%%%%%%%%%%%%%%%%%%%%%%%%
% Slide 3: Design Optimization using Genetic Algorithms
%%%%%%%%%%%%%%%%%%%%%%%%%%%%%%%%%%%%%%%%%%%%%%%%%%%%%%%%%%%
\begin{frame}[fragile]\frametitle{Design Optimization using Genetic Algorithms}
    \begin{itemize}
        \item \textbf{Problem}: a design objective to minimize, for example the weight of a beam under stress and deflection limits
        \item \textbf{Source}: \href{https://deap.readthedocs.io/}{DEAP documentation}
        \item \textbf{Description}: Optimize engineering designs using genetic algorithms to evaluate performance metrics.
    \end{itemize}
    \begin{lstlisting}[language=Python]
from deap import base, creator, tools, algorithms

creator.create("FitnessMin", base.Fitness, weights=(-1.0,))
creator.create("Individual", list, fitness=creator.FitnessMin)

toolbox = base.Toolbox()
# Define the design variables and the objective (fitness), then register
# individual, population, evaluate, mate, mutate and select on the toolbox
population = toolbox.population(n=50)
algorithms.eaSimple(population, toolbox, cxpb=0.5, mutpb=0.2, ngen=40)
    \end{lstlisting}
\end{frame}

%%%%%%%%%%%%%%%%%%%%%%%%%%%%%%%%%%%%%%%%%%%%%%%%%%%%%%%%%%%
% Slide 4: Fault Detection in Rotating Machinery using SVM
%%%%%%%%%%%%%%%%%%%%%%%%%%%%%%%%%%%%%%%%%%%%%%%%%%%%%%%%%%%
\begin{frame}[fragile]\frametitle{Fault Detection in Rotating Machinery using SVM}
    \begin{itemize}
        \item \textbf{Dataset}: Case Western Reserve University Bearing Data (vibration signals from healthy and faulty bearings)
        \item \textbf{Source}: \href{https://engineering.case.edu/bearingdatacenter}{CWRU Bearing Data Center}
        \item \textbf{Description}: Classify faults in rotating machinery using support vector machines (SVM).
    \end{itemize}
    \begin{lstlisting}[language=Python]
import pandas as pd
from sklearn.svm import SVC
from sklearn.model_selection import train_test_split

data = pd.read_csv('machinery_fault_data.csv')
X = data.drop('fault_type', axis=1)
y = data['fault_type']

X_train, X_test, y_train, y_test = train_test_split(X, y, test_size=0.25)
model = SVC()
model.fit(X_train, y_train)
accuracy = model.score(X_test, y_test)
    \end{lstlisting}
\end{frame}

%%%%%%%%%%%%%%%%%%%%%%%%%%%%%%%%%%%%%%%%%%%%%%%%%%%%%%%%%%%
% Slide 5: Thermal Analysis using Neural Networks
%%%%%%%%%%%%%%%%%%%%%%%%%%%%%%%%%%%%%%%%%%%%%%%%%%%%%%%%%%%
\begin{frame}[fragile]\frametitle{Thermal Analysis using Neural Networks}
    \begin{itemize}
        \item \textbf{Dataset}: A public thermal conductivity dataset
        \item \textbf{Source}: \href{https://www.kaggle.com/datasets?search=thermal+conductivity}{Kaggle search}
        \item \textbf{Description}: Predict thermal conductivity of materials using neural networks.
    \end{itemize}
    \begin{lstlisting}[language=Python]
from keras.models import Sequential
from keras.layers import Dense

model = Sequential()
model.add(Dense(64, activation='relu', input_dim=10))
model.add(Dense(32, activation='relu'))
model.add(Dense(1))

model.compile(optimizer='adam', loss='mean_squared_error')
    \end{lstlisting}
\end{frame}

%%%%%%%%%%%%%%%%%%%%%%%%%%%%%%%%%%%%%%%%%%%%%%%%%%%%%%%%%%%
% Slide 6: Quality Control in Manufacturing using CNNs
%%%%%%%%%%%%%%%%%%%%%%%%%%%%%%%%%%%%%%%%%%%%%%%%%%%%%%%%%%%
\begin{frame}[fragile]\frametitle{Quality Control in Manufacturing using CNNs}
    \begin{itemize}
        \item \textbf{Dataset}: A public manufacturing defect image dataset
        \item \textbf{Source}: \href{https://www.kaggle.com/datasets?search=defect+detection}{Kaggle search}
        \item \textbf{Description}: Train a CNN to classify manufacturing defects based on part images.
    \end{itemize}
    \begin{lstlisting}[language=Python]
from keras.models import Sequential
from keras.layers import Conv2D, MaxPooling2D, Flatten, Dense

model = Sequential()
model.add(Conv2D(32, (3,3), activation='relu', input_shape=(128,128,3)))
model.add(MaxPooling2D(pool_size=(2,2)))
model.add(Flatten())
model.add(Dense(128, activation='relu'))
model.add(Dense(10, activation='softmax'))   # 10 = number of defect classes

model.compile(optimizer='adam', loss='categorical_crossentropy', metrics=['accuracy'])
    \end{lstlisting}
\end{frame}

%%%%%%%%%%%%%%%%%%%%%%%%%%%%%%%%%%%%%%%%%%%%%%%%%%%%%%%%%%%
% Slide 7: Energy Consumption Prediction using Time Series Analysis
%%%%%%%%%%%%%%%%%%%%%%%%%%%%%%%%%%%%%%%%%%%%%%%%%%%%%%%%%%%
\begin{frame}[fragile]\frametitle{Energy Consumption Prediction using Time Series Analysis}
    \begin{itemize}
        \item \textbf{Dataset}: Individual Household Electric Power Consumption
        \item \textbf{Source}: \href{https://archive.ics.uci.edu/dataset/235/individual+household+electric+power+consumption}{UCI Machine Learning Repository}
        \item \textbf{Description}: Forecast energy consumption patterns using ARIMA models.
    \end{itemize}
    \begin{lstlisting}[language=Python]
import pandas as pd
from statsmodels.tsa.arima.model import ARIMA

data = pd.read_csv('energy_consumption.csv')
model = ARIMA(data['consumption'], order=(5,1,0))
model_fit = model.fit()
forecast = model_fit.forecast(steps=10)
    \end{lstlisting}
\end{frame}

%%%%%%%%%%%%%%%%%%%%%%%%%%%%%%%%%%%%%%%%%%%%%%%%%%%%%%%%%%%
% Slide 8: Material Property Prediction using Gaussian Processes
%%%%%%%%%%%%%%%%%%%%%%%%%%%%%%%%%%%%%%%%%%%%%%%%%%%%%%%%%%%
\begin{frame}[fragile]\frametitle{Material Property Prediction using Gaussian Processes}
    \begin{itemize}
        \item \textbf{Dataset}: Materials Project Database
        \item \textbf{Source}: \href{https://materialsproject.org/}{Materials Project}
        \item \textbf{Description}: Predict material properties using Gaussian process regression.
    \end{itemize}
    \begin{lstlisting}[language=Python]
from sklearn.gaussian_process import GaussianProcessRegressor

# X_train, y_train, X_test: loaded from the Materials Project dataset above
gp = GaussianProcessRegressor()
gp.fit(X_train, y_train)
predictions = gp.predict(X_test)
    \end{lstlisting}
\end{frame}

%%%%%%%%%%%%%%%%%%%%%%%%%%%%%%%%%%%%%%%%%%%%%%%%%%%%%%%%%%%
% Slide 9: Load Prediction in Structural Engineering using Random Forests
%%%%%%%%%%%%%%%%%%%%%%%%%%%%%%%%%%%%%%%%%%%%%%%%%%%%%%%%%%%
\begin{frame}[fragile]\frametitle{Strength Prediction in Structural Engineering using Random Forests}
    \begin{itemize}
        \item \textbf{Dataset}: Concrete Compressive Strength
        \item \textbf{Source}: \href{https://archive.ics.uci.edu/dataset/165/concrete+compressive+strength}{UCI Machine Learning Repository}
        \item \textbf{Description}: Predict the compressive strength of concrete from its mix using random forests.
    \end{itemize}
    \begin{lstlisting}[language=Python]
import pandas as pd
from sklearn.ensemble import RandomForestRegressor
from sklearn.model_selection import train_test_split

data = pd.read_csv('concrete_data.csv')
X = data.drop('strength', axis=1)
y = data['strength']

X_train, X_test, y_train, y_test = train_test_split(X, y, test_size=0.2)
model = RandomForestRegressor()
model.fit(X_train, y_train)
predictions = model.predict(X_test)
    \end{lstlisting}
\end{frame}

%%%%%%%%%%%%%%%%%%%%%%%%%%%%%%%%%%%%%%%%%%%%%%%%%%%%%%%%%%%
% Slide 10: Simulation of Fluid Dynamics using Deep Learning
%%%%%%%%%%%%%%%%%%%%%%%%%%%%%%%%%%%%%%%%%%%%%%%%%%%%%%%%%%%
\begin{frame}[fragile]\frametitle{Simulation of Fluid Dynamics using Deep Learning}
    \begin{itemize}
        \item \textbf{Dataset}: Fluid Dynamics Simulation Data
        \item \textbf{Source}: \href{https://www.openfoam.com/}{OpenFOAM}
        \item \textbf{Description}: Simulate fluid dynamics using deep learning models.
    \end{itemize}
    \begin{lstlisting}[language=Python]
import tensorflow as tf

model = tf.keras.Sequential([
    tf.keras.layers.Conv2D(32, (3,3), activation='relu', input_shape=(64,64,1)),
    tf.keras.layers.Flatten(),
    tf.keras.layers.Dense(64, activation='relu'),
    tf.keras.layers.Dense(1)
])

model.compile(optimizer='adam', loss='mean_squared_error')
    \end{lstlisting}
\end{frame}

%%%%%%%%%%%%%%%%%%%%%%%%%%%%%%%%%%%%%%%%%%%%%%%%%%%%%%%%%%%
\begin{frame}[fragile]\frametitle{Mech References}
    \begin{enumerate}
        \item \href{https://github.com/hanhuark/MEEG-54403}{GitHub: MEEG-54403 Course Materials}
        \item \href{https://www.neuralconcept.com/post/applications-of-machine-learning-in-mechanical-engineering}{Neural Concept: Applications of ML in Mechanical Engineering}
        \item \href{https://msse.gatech.edu/me8813ML/index.html}{Georgia Tech: ME 8813 Machine Learning for Engineers}
        \item \href{https://lms.decibelslab.com/courses/Machine-Learning-for-Mechanical-Engineers}{Decibels Lab: ML for Mechanical Engineers}
        \item \href{https://www.meche.engineering.cmu.edu/research/machine-learning.html}{Carnegie Mellon: ML in Mechanical Engineering Research}
        \item \href{https://www.wiredworkers.io/blog/machine-learning-in-mechanical-engineering/}{Wired Workers: Blog on ML in Mechanical Engineering}
        \item \href{https://github.com/Parnni/Projects}{GitHub: Projects on Machine Learning}
        \item \href{https://stanford-me343.github.io}{Stanford: ME 343 Advanced ML Applications}
    \end{enumerate}
\end{frame}


% Agriculture assignments and references (12 frames): restored for the final session
% (agriculture is in scope alongside mechanical and manufacturing).
%%%%%%%%%%%%%%%%%%%%%%%%%%%%%%%%%%%%%%%%%%%%%%%%%%%%%%%%%%%
\begin{frame}[fragile]\frametitle{}
\begin{center}
{\Large Assignments in Agriculture}
\end{center}
\end{frame}
 
%%%%%%%%%%%%%%%%%%%%%%%%%%%%%%%%%%%%%%%%%%%%%%%%%%%%%%%%%%%
% Slide 1: Crop Yield Prediction using Linear Regression
%%%%%%%%%%%%%%%%%%%%%%%%%%%%%%%%%%%%%%%%%%%%%%%%%%%%%%%%%%%
\begin{frame}[fragile]\frametitle{Crop Yield Prediction using Linear Regression}
    \begin{itemize}
        \item \textbf{Dataset}: Crop Yield Prediction Dataset
        \item \textbf{Source}: \href{https://www.kaggle.com/datasets?search=crop+yield+prediction}{Kaggle search}
        \item \textbf{Description}: Predict crop yields based on environmental factors such as temperature, humidity, and rainfall using linear regression.
    \end{itemize}
    \begin{lstlisting}[language=Python]
import pandas as pd
from sklearn.model_selection import train_test_split
from sklearn.linear_model import LinearRegression

data = pd.read_csv('crop_yield.csv')
X = data[['temperature', 'humidity', 'rainfall']]
y = data['yield']

X_train, X_test, y_train, y_test = train_test_split(X, y, test_size=0.2)
model = LinearRegression()
model.fit(X_train, y_train)
predictions = model.predict(X_test)
    \end{lstlisting}
\end{frame}

%%%%%%%%%%%%%%%%%%%%%%%%%%%%%%%%%%%%%%%%%%%%%%%%%%%%%%%%%%%
% Slide 2: Disease Detection in Crops using Decision Trees
%%%%%%%%%%%%%%%%%%%%%%%%%%%%%%%%%%%%%%%%%%%%%%%%%%%%%%%%%%%
\begin{frame}[fragile]\frametitle{Disease Detection in Crops using Decision Trees}
    \begin{itemize}
        \item \textbf{Dataset}: Plant Village Dataset
        \item \textbf{Source}: \href{https://www.kaggle.com/datasets/emmarex/plantdisease}{Kaggle}
        \item \textbf{Description}: Use decision trees to classify crop diseases based on features extracted from images of leaves.
    \end{itemize}
    \begin{lstlisting}[language=Python]
from sklearn.tree import DecisionTreeClassifier
from sklearn.model_selection import train_test_split
import pandas as pd

data = pd.read_csv('plant_disease.csv')
X = data.drop('disease', axis=1)
y = data['disease']

X_train, X_test, y_train, y_test = train_test_split(X, y, test_size=0.2)
model = DecisionTreeClassifier()
model.fit(X_train, y_train)
accuracy = model.score(X_test, y_test)
    \end{lstlisting}
\end{frame}

%%%%%%%%%%%%%%%%%%%%%%%%%%%%%%%%%%%%%%%%%%%%%%%%%%%%%%%%%%%
% Slide 3: Weed Detection using Convolutional Neural Networks (CNN)
%%%%%%%%%%%%%%%%%%%%%%%%%%%%%%%%%%%%%%%%%%%%%%%%%%%%%%%%%%%
\begin{frame}[fragile]\frametitle{Weed Detection using Convolutional Neural Networks (CNN)}
    \begin{itemize}
        \item \textbf{Dataset}: Weeds Dataset
        \item \textbf{Source}: \href{https://www.kaggle.com/datasets?search=weed+detection}{Kaggle search}
        \item \textbf{Description}: Implement a CNN to detect weeds in crop fields by training on images of weeds versus crops.
    \end{itemize}
    \begin{lstlisting}[language=Python]
from keras.models import Sequential
from keras.layers import Conv2D, MaxPooling2D, Flatten, Dense

model = Sequential()
model.add(Conv2D(32, (3,3), activation='relu', input_shape=(64,64,3)))
model.add(MaxPooling2D(pool_size=(2,2)))

model.add(Flatten())
model.add(Dense(units=128, activation='relu'))
model.add(Dense(units=1, activation='sigmoid'))

model.compile(optimizer='adam', loss='binary_crossentropy', metrics=['accuracy'])
    \end{lstlisting}
\end{frame}

%%%%%%%%%%%%%%%%%%%%%%%%%%%%%%%%%%%%%%%%%%%%%%%%%%%%%%%%%%%
% Slide 4: Soil Quality Prediction using Random Forests
%%%%%%%%%%%%%%%%%%%%%%%%%%%%%%%%%%%%%%%%%%%%%%%%%%%%%%%%%%%
\begin{frame}[fragile]\frametitle{Soil Quality Prediction using Random Forests}
    \begin{itemize}
        \item \textbf{Dataset}: Soil Quality Dataset
        \item \textbf{Source}: \href{https://www.kaggle.com/datasets?search=soil+quality}{Kaggle search}
        \item \textbf{Description}: Use random forests to predict soil quality based on various chemical and physical properties of the soil.
    \end{itemize}
    \begin{lstlisting}[language=Python]
import pandas as pd
from sklearn.ensemble import RandomForestClassifier
from sklearn.model_selection import train_test_split

data = pd.read_csv('soil_quality.csv')
X = data.drop('quality', axis=1)
y = data['quality']

X_train, X_test, y_train, y_test = train_test_split(X, y, test_size=0.3)
model = RandomForestClassifier()
model.fit(X_train, y_train)

accuracy = model.score(X_test, y_test)
    \end{lstlisting}
\end{frame}

%%%%%%%%%%%%%%%%%%%%%%%%%%%%%%%%%%%%%%%%%%%%%%%%%%%%%%%%%%%
% Slide 5: Irrigation Scheduling using Support Vector Machines (SVM)
%%%%%%%%%%%%%%%%%%%%%%%%%%%%%%%%%%%%%%%%%%%%%%%%%%%%%%%%%%%
\begin{frame}[fragile]\frametitle{Irrigation Scheduling using Support Vector Machines (SVM)}
    \begin{itemize}
        \item \textbf{Dataset}: Irrigation Scheduling Dataset
        \item \textbf{Source}: \href{https://www.kaggle.com/datasets?search=irrigation+scheduling}{Kaggle search}
        \item \textbf{Description}: Implement SVM to determine optimal irrigation schedules based on weather and soil moisture data.
    \end{itemize}
    \begin{lstlisting}[language=Python]
from sklearn.svm import SVC
from sklearn.model_selection import train_test_split
import pandas as pd

data = pd.read_csv('irrigation_data.csv')
X = data.drop('schedule', axis=1)
y = data['schedule']

X_train, X_test, y_train, y_test = train_test_split(X, y, test_size=0.2)
model = SVC()
model.fit(X_train, y_train)

accuracy = model.score(X_test, y_test)
    \end{lstlisting}
\end{frame}

%%%%%%%%%%%%%%%%%%%%%%%%%%%%%%%%%%%%%%%%%%%%%%%%%%%%%%%%%%%
% Slide 6: Crop Disease Prediction using K-Nearest Neighbors (KNN)
%%%%%%%%%%%%%%%%%%%%%%%%%%%%%%%%%%%%%%%%%%%%%%%%%%%%%%%%%%%
\begin{frame}[fragile]\frametitle{Crop Disease Prediction using K-Nearest Neighbors (KNN)}
    \begin{itemize}
        \item \textbf{Dataset}: Crop Disease Dataset
        \item \textbf{Source}: \href{https://www.kaggle.com/datasets?search=crop+disease+prediction}{Kaggle search}
        \item \textbf{Description}: Use KNN to classify different diseases affecting crops based on features like leaf color and texture.
    \end{itemize}
    \begin{lstlisting}[language=Python]
from sklearn.neighbors import KNeighborsClassifier
from sklearn.model_selection import train_test_split
import pandas as pd

data = pd.read_csv('crop_disease_data.csv')
X = data.drop('disease', axis=1)
y = data['disease']

X_train, X_test, y_train, y_test = train_test_split(X, y, test_size=0.25)
model = KNeighborsClassifier(n_neighbors=5)
model.fit(X_train, y_train)

accuracy = model.score(X_test, y_test)
    \end{lstlisting}
\end{frame}

%%%%%%%%%%%%%%%%%%%%%%%%%%%%%%%%%%%%%%%%%%%%%%%%%%%%%%%%%%%
% Slide 7: Yield Prediction using Gradient Boosting Machines (GBM)
%%%%%%%%%%%%%%%%%%%%%%%%%%%%%%%%%%%%%%%%%%%%%%%%%%%%%%%%%%%
\begin{frame}[fragile]\frametitle{Yield Prediction using Gradient Boosting Machines (GBM)}
    \begin{itemize}
        \item \textbf{Dataset}: US Crop Yield Data
        \item \textbf{Source}: \href{https://www.nass.usda.gov/}{USDA}
        \item \textbf{Description}: Apply gradient boosting to predict crop yields based on historical yield data and environmental factors.
    \end{itemize}
    \begin{lstlisting}[language=Python]
import pandas as pd
from sklearn.ensemble import GradientBoostingRegressor
from sklearn.model_selection import train_test_split

data = pd.read_csv('us_crop_yield.csv')
X = data.drop('yield', axis=1)
y = data['yield']

X_train, X_test, y_train, y_test = train_test_split(X, y, test_size=0.2)
model = GradientBoostingRegressor()
model.fit(X_train, y_train)

predictions = model.predict(X_test)
    \end{lstlisting}
\end{frame}

%%%%%%%%%%%%%%%%%%%%%%%%%%%%%%%%%%%%%%%%%%%%%%%%%%%%%%%%%%%
% Slide 8: Pest Detection using Image Classification with CNNs
%%%%%%%%%%%%%%%%%%%%%%%%%%%%%%%%%%%%%%%%%%%%%%%%%%%%%%%%%%%
\begin{frame}[fragile]\frametitle{Pest Detection using Image Classification with CNNs}
    \begin{itemize}
        \item \textbf{Dataset}: Pest Image Dataset
        \item \textbf{Source}: \href{https://www.kaggle.com/datasets?search=pest+detection}{Kaggle search}
        \item \textbf{Description}: Train a CNN for pest detection in crops by classifying images of healthy versus infested plants.
    \end{itemize}
    \begin{lstlisting}[language=Python]
from keras.models import Sequential
from keras.layers import Conv2D, MaxPooling2D, Flatten, Dense

model = Sequential()
model.add(Conv2D(32,(3,3), activation='relu', input_shape=(128,128,3)))
model.add(MaxPooling2D(pool_size=(2,2)))
model.add(Flatten())
model.add(Dense(units=128 , activation='relu'))
model.add(Dense(units=1, activation='sigmoid'))   # healthy vs infested

model.compile(optimizer='adam', loss='binary_crossentropy', metrics=['accuracy'])
    \end{lstlisting}
\end{frame}

%%%%%%%%%%%%%%%%%%%%%%%%%%%%%%%%%%%%%%%%%%%%%%%%%%%%%%%%%%%
% Slide 9: Precision Agriculture using Clustering Techniques (K-Means)
%%%%%%%%%%%%%%%%%%%%%%%%%%%%%%%%%%%%%%%%%%%%%%%%%%%%%%%%%%%
\begin{frame}[fragile]\frametitle{Precision Agriculture using Clustering Techniques (K-Means)}
    \begin{itemize}
        \item \textbf{Dataset}: Precision Agriculture Dataset
        \item \textbf{Source}: \href{https://www.kaggle.com/datasets?search=precision+agriculture}{Kaggle search}
        \item \textbf{Description}: Use K-means clustering to segment agricultural land into different zones for precision farming practices.
    \end{itemize}
    \begin{lstlisting}[language=Python]
import pandas as pd
from sklearn.cluster import KMeans

data = pd.read_csv('precision_agriculture.csv')
X = data[['latitude', 'longitude']]

kmeans = KMeans(n_clusters=3) # Number of clusters can be adjusted.
kmeans.fit(X)

data['cluster'] = kmeans.labels_
    \end{lstlisting}
\end{frame}

%%%%%%%%%%%%%%%%%%%%%%%%%%%%%%%%%%%%%%%%%%%%%%%%%%%%%%%%%%%
% Slide 10: Weather Forecasting for Agriculture using LSTM Networks
%%%%%%%%%%%%%%%%%%%%%%%%%%%%%%%%%%%%%%%%%%%%%%%%%%%%%%%%%%%
\begin{frame}[fragile]\frametitle{Weather Forecasting for Agriculture using LSTM Networks}
    \begin{itemize}
        \item \textbf{Dataset}: Weather Data
        \item \textbf{Source}: \href{https://www.ncdc.noaa.gov/cdo-web/}{NOAA}
        \item \textbf{Description}: Implement LSTM networks to forecast weather conditions that affect agricultural practices like planting and harvesting.
    \end{itemize}
    \begin{lstlisting}[language=Python]
import numpy as np
import pandas as pd
from keras.models import Sequential
from keras.layers import LSTM,Dense

data = pd.read_csv('weather_data.csv')
X, y = data.iloc[:, :-1].values, data.iloc[:, -1].values
X = X.reshape((X.shape[0], 1, X.shape[1]))   # LSTM input: (samples, timesteps, features)

model = Sequential()
model.add(LSTM(50, input_shape=(X.shape[1], X.shape[2])))
model.add(Dense(1))
model.compile(loss='mean_squared_error', optimizer='adam')
    \end{lstlisting}
\end{frame}


%%%%%%%%%%%%%%%%%%%%%%%%%%%%%%%%%%%%%%%%%%%%%%%%%%%%%%%%%%%
\begin{frame}[fragile]\frametitle{References}
    \begin{enumerate}
        \item \href{https://www.itransition.com/machine-learning/agriculture}{Itransition: Machine Learning in Agriculture}
        \item \href{https://www.simplilearn.com/10-algorithms-machine-learning-engineers-need-to-know-article}{Simplilearn: Key ML Algorithms}
        \item \href{https://www.ijarcce.com/upload/2016/april-16/IJARCCE%2087.pdf}{IJARCCE: Research Paper on ML Applications in Agriculture}
        \item \href{https://www.datacamp.com/blog/top-machine-learning-use-cases-and-algorithms}{DataCamp: Top ML Use Cases}
        \item \href{https://indatalabs.com/blog/ml-in-agriculture}{InData Labs: ML in Agriculture}
        \item \href{https://www.projectpro.io/article/top-10-machine-learning-projects-for-beginners-in-2021/397}{ProjectPro: Top ML Projects}
        \item \href{https://www.croptracker.com/blog/what-is-machine-learning-how-will-it-benefit-agriculture.html}{CropTracker: ML Benefits in Agriculture}
        \item \href{https://www.toptal.com/machine-learning/machine-learning-theory-an-introductory-primer}{Toptal: Introduction to ML Theory}
		\item \href{https://www.youtube.com/watch?v=zcq5aw9t-Ds}{Plant Leaf Disease Detection Using CNN | Python - YouTube}		
		\item \href{https://www.youtube.com/watch?v=T7\_GR8l-N5A}{Project 20 : Crop Recommendation Using Machine Learning - YouTube}		
		\item \href{https://www.youtube.com/watch?v=LEnjAR6AwOo}{Project 23: Crop Yield Prediction Using Machine Learning | KNOWLEDGE DOCTOR - YouTube}
		\item \href{https://www.youtube.com/watch?v=Yj-lHSfovw4}{Deep Learning Project:Cotton Disease Prediction Using Pytorch l Web Apps in Flask | KNOWLEDGE DOCTOR - YouTube}
		\item \href{https://www.youtube.com/watch?v=59H5\_1Q4lqU}{IoT Based Smart Plant Monitoring System | Arduino UNO | Bangla & English Version | KNOWLEDGE DOCTOR - YouTube}
    \end{enumerate}
\end{frame}


% Quick check after the assignments (upgrade-deck Task 6)
\begin{frame}[fragile]\frametitle{Quick Check: Assignments}
\begin{block}{Think About It}
You build a crop-yield model from ten years of farm records, and a random split gives a very high score.
Which split best tests whether the model will work next season?
\begin{enumerate}
  \item[A.] Keep the random split, but repeat it ten times.
  \item[B.] Train on the earlier years and test on the latest year or two.
  \item[C.] Test on the same years the model was trained on, to check it learned them.
  \item[D.] Shuffle the rows again until the score is stable.
\end{enumerate}
\end{block}
\end{frame}

\begin{frame}[fragile]\frametitle{Quick Check: Assignments}
\begin{block}{Answer}
\textbf{Correct: B.} A forecast is always made for a season the model has not seen, so the test set must lie
later in time. Random splits (A, D) let neighbouring seasons leak across the split and flatter the score, and
testing on the training years (C) measures memory, not prediction.
\end{block}
\end{frame}

% Closing: recap, career paths, next steps (last day of the series)
\section[Career]{Where to Go From Here}
% Closing material for the last MLCoEP session: recap, careers, next steps.
% No salary figures or market statistics are quoted: they date quickly.

%%%%%%%%%%%%%%%%%%%%%%%%%%%%%%%%%%%%%%%%%%%%%%%%%%%%%%%%%%%
\begin{frame}[fragile]\frametitle{}
\begin{center}
{\Large Where to Go From Here}

\medskip
{\small (What you can do now, and how to keep growing)}
\end{center}
\end{frame}

%%%%%%%%%%%%%%%%%%%%%%%%%%%%%%%%%%%%%%%%%%%%%%%%%%%
\begin{frame}[fragile]\frametitle{What You Can Now Do}
\begin{itemize}
\item Frame an engineering question as a learning problem: predict a number, assign a class, or find groups.
\item Pick a sensible model, train it, and check it on data it has not seen.
\item Read a result critically: over-fitting, leakage, and wrong metrics.
\item Connect sensors, machines and fields to data, from IoT streams to crop records.
\end{itemize}
\end{frame}

%%%%%%%%%%%%%%%%%%%%%%%%%%%%%%%%%%%%%%%%%%%%%%%%%%%
\begin{frame}[fragile]\frametitle{The Same Workflow, Three Settings}
\begin{center}
\resizebox{\linewidth}{!}{%
\begin{tabular}{llll}
\toprule
Step & Mechanical & Manufacturing & Agriculture \\
\midrule
Question & Will this bearing fail? & Is this part defective? & What yield to expect? \\
Data & Vibration, temperature & Line sensors, images & Weather, soil, satellite \\
Model & Regression, SVM & CNN, anomaly detection & Regression, random forest \\
Check & Time-split test set & Defect-rate cost & Season-wise hold-out \\
\bottomrule
\end{tabular}}
\end{center}
\end{frame}

%%%%%%%%%%%%%%%%%%%%%%%%%%%%%%%%%%%%%%%%%%%%%%%%%%%
\begin{frame}[fragile]\frametitle{Career Paths That Use This}
\begin{itemize}
\item \textbf{Reliability and predictive maintenance:} models that warn before a machine fails.
\item \textbf{Quality and process engineering:} defect detection, process tuning, yield improvement.
\item \textbf{Digital twins and simulation:} learned surrogates that speed up design studies.
\item \textbf{Agritech and precision farming:} yield, irrigation, disease and pest models.
\item \textbf{Data and ML engineering:} building and running the data pipelines behind these systems.
\end{itemize}
\end{frame}

%%%%%%%%%%%%%%%%%%%%%%%%%%%%%%%%%%%%%%%%%%%%%%%%%%%
\begin{frame}[fragile]\frametitle{Advice: Be the Engineer Who Also Knows ML}
\begin{itemize}
\item Domain depth is your edge: a data scientist cannot tell a worn tool from a bad sensor, and you can.
\item Keep your engineering fundamentals sharp: thermodynamics, mechanics, controls and materials still decide what the model should predict.
\item Learn to say what a model cannot do. Honest limits build trust faster than bold claims.
\item Start small, in a real plant or on a real farm: one problem, one dataset, one measurable result.
\end{itemize}
\end{frame}

%%%%%%%%%%%%%%%%%%%%%%%%%%%%%%%%%%%%%%%%%%%%%%%%%%%
\begin{frame}[fragile]\frametitle{Skills to Keep Building}
\begin{itemize}
\item \textbf{Python and data handling:} pandas, NumPy, plotting, reading messy sensor files.
\item \textbf{Statistics:} sampling, uncertainty, hypothesis tests, and how to design an experiment.
\item \textbf{Modelling:} scikit-learn first, then a deep learning library once a problem needs it.
\item \textbf{Deployment basics:} saving a model, serving it, and watching it over time.
\item \textbf{Communication:} one clear plot and one clear sentence for a non-specialist manager.
\end{itemize}
\end{frame}

%%%%%%%%%%%%%%%%%%%%%%%%%%%%%%%%%%%%%%%%%%%%%%%%%%%
\begin{frame}[fragile]\frametitle{A Portfolio That Shows Your Work}
\begin{itemize}
\item Pick two or three small projects from the assignment ideas, and finish them properly.
\item For each: state the problem, the data source, the method, the result, and what went wrong.
\item Put the code and a short write-up in a public repository, with setup notes that a stranger can follow.
\item Ideas: tool wear estimation, bearing fault detection, crop yield prediction, defect classification from images.
\end{itemize}
\end{frame}

%%%%%%%%%%%%%%%%%%%%%%%%%%%%%%%%%%%%%%%%%%%%%%%%%%%
\begin{frame}[fragile]\frametitle{How to Keep Learning}
\begin{itemize}
\item Work with public datasets: Kaggle, UCI, NASA prognostics repository.
\item Read one short paper or case study a week from your own area, and note what data it used.
\item Rebuild a published result before trying to improve on it.
\item Join a community: a local meetup, a college club or an open-source project.
\item Revisit the worked examples: they read differently once the full workflow is clear.
\end{itemize}
\end{frame}

%%%%%%%%%%%%%%%%%%%%%%%%%%%%%%%%%%%%%%%%%%%%%%%%%%%
\begin{frame}[fragile]\frametitle{Responsible Practice}
\begin{itemize}
\item \textbf{Safety first:} a model that controls or advises on a machine needs a safe fallback when it is wrong.
\item \textbf{Data care:} sensor and farm data can be commercially sensitive; get permission before sharing.
\item \textbf{Fairness to people:} models that rank workers or farmers can disadvantage groups; ask who is affected.
\item \textbf{Document:} data source, assumptions and limits go with every model you hand over.
\end{itemize}
\end{frame}

%%%%%%%%%%%%%%%%%%%%%%%%%%%%%%%%%%%%%%%%%%%%%%%%%%%
\begin{frame}[fragile]\frametitle{Quick Check: Your First Project}
\begin{block}{Think About It}
You have one month and one small dataset of bearing vibration readings from a local workshop.
What would you do first, and what would you report to show that your model can be trusted?
\end{block}
\end{frame}

%%%%%%%%%%%%%%%%%%%%%%%%%%%%%%%%%%%%%%%%%%%%%%%%%%%
\begin{frame}[fragile]\frametitle{Quick Check: Your First Project}
\begin{block}{Answer}
Look at the data first: plot the signals, check for missing values and mislabelled runs.
Start with a simple baseline, such as a few engineered features and a regularised model,
before trying anything deep. Split by machine or by time, not at random, so the test data
are truly new. Report the metric that matches the cost of a missed failure, and show the
cases where the model is wrong, so that its limits are visible.
\end{block}
\end{frame}

%%%%%%%%%%%%%%%%%%%%%%%%%%%%%%%%%%%%%%%%%%%%%%%%%%%
\begin{frame}[fragile]\frametitle{Thank You, and Good Luck}
\begin{center}
{\Large You now have the tools; the problems are waiting.}
\end{center}
\begin{itemize}
\item Take one problem from your own surroundings and start this week.
\item Share what you build: a failed attempt, honestly written up, is still useful to others.
\end{itemize}
\end{frame}

